\subsection{一交換子の平方追加Segre射影}
\begin{proposition}\label{prop:cubic-CY-one-segre-square-empty}
一交換子を持つCY包絡で、
$(U_0,U)=(\langle ac,ad,bc,bd\rangle,\langle ac,ad,bc,bd,a^2\rangle)$
となるものはH3を与えない。
\end{proposition}
\begin{proof}
固定出典\cite{AS4CYpinned}の平方追加Segreの全射影被覆の必要方向を用いる。
横断枝は $A_h,N,B_1,B_2$、内部枝は $S_a,S_{a+b}$ と七孤立表示で尽くされる。

元六関係から三次・四次の全関係挿入行列を作り、商への座標射影と
中心条件の行列を計算する。パラメータ族では有理関数体上の商射影を、
その分母が消えない範囲で特殊化できる。
各射影は元関係を消し、選んだ自由列上で恒等写像である。
従って実際のAS4点では次数別次元20,35から商との同型が従う。
中心条件の非零20次小行列式の零点と、商射影の分母の零点を全て取り、
許容値を元関係から再計算した。各既約因子はその数体上で検査するので全共役根を含む。
検証器\\
\texttt{verify\_cubic\_CY\_one\_segre\_square.py}\quad
に係数表示と再現手順を与えた。

横断枝と七孤立表示では全値で中心三次が零である。
内部二族でも $48h^2+1=0$ 以外に中心三次はない。
ここで $S_{a+b}$ の中心三次は $\langle a^3,b^3\rangle$。
任意の非零二元三次式は、可換な一次元生成元 $a,b$ の張る平面上で積に分解する。
三次切断後も一次・二次は変わらないため、整域にはならない。

残る $S_a$ の中心三次は $\langle a^3,b^3,c^3,d^3\rangle$ である。
$q=(4h+1)/(4h-1)$ は原始三乗根で、固定出典の(I1)から
\[
 ca=qac,\quad cb=q^{-1}bc,\quad da=q^{-1}ad,\quad db=qbd,
 \qquad [a,b]=0,\quad dc=cd+\kappa a^2,\qquad \kappa\ne0
\]
を得る。$G=Aa^3+Bb^3+Cc^3+Dd^3$ の切断が整域であると仮定する。
$C=D=0$ は上と同じ積分解に矛盾する。
$c,d$ を必要なら交換して $D\ne0$ とする。この交換は $q$ と $\kappa$ の値を変更するだけである。
非零正常元 $a$ を逆元化し、次数零部分で
\[
 x=ca^{-1},\quad y=da^{-1},\quad z=ba^{-1}
\]
とおく。すると $xz=qzx$, $yz=q^{-1}zy$ および
$yx=q^2xy+q\kappa$ となる。
$L=q\kappa/(1-q^2)$、$u=xy-L$ とおけば、
$ux=q^2xu$, $uy=qyu$, $uz=zu$ である。
更に非零中央元 $z^3$ を逆元化すると $v=uz^{-1}$ は
$x,y,z$ と可換である。
$X=x^3,Y=y^3,Z=z^3$ と書くと
\[
 XY=L^3+v^3 Z,\qquad A+BZ+CX+DY=0.
\]
第一式は $x^3 y^3$ を $xy$ の三つの共役の積として得る。
第二式は $Ga^{-3}=0$ である。

$x,y,z$ の生成する局所化部分環では、$X,Z,v$ の生成する可換部分環は中央にある。
その分数体を $K$ とし、中央局所化しても整域であると仮定したことに注意する。
$X$ は $k$ 上超越的である。もし代数的なら $X=\alpha^3$ となり、
量子二項式により $(c-\alpha a)^3=c^3-\alpha^3a^3=0$ となる。
これは一次の独立性と整域性に反する。
従って
\[
 Z=\frac{-L^3-(A/D)X-(C/D)X^2}{v^3+(B/D)X}
\]
であり、分母の非零性は左辺との積が非零多項式であることから従う。
$X$ は $K$ で立方ではない。立方なら $x^3=X$ の因数分解により $x\in K$ となり、
非零の $z$ との交換式に反する。
従って $K(x)/K$ は三次巡回拡大で、$z$ の共役作用がその生成元である。
分子は $k$ 上で $X$ の一次因子に分解し、全て $K(x)/K$ のノルムとなる。
実際 $N(x-\alpha)=X-\alpha^3$ であり、定数も代数閉性によりノルムである。
分母も $N(v+\beta x)$、$\beta^3=B/D$ である。
従って $Z=N(w)$ となる $w\in K(x)^*$ が存在する。
すると $(w^{-1}z)^3=1$ であり、整域性から $w^{-1}z$ はスカラーでなければならない。
これは $z$ が $x$ に非自明な共役作用を持つことに矛盾する。
従って最後の例外でも整域切断は存在しない。
\end{proof}
